\documentclass{sajzk}

\usepackage{array}
\usepackage{mathtools}
\usepackage{minted}

\usetikzlibrary {arrows.meta, angles, quotes}
\usetikzlibrary {calc}

\begin{document}
\section{Basisvektoren für GA in 3D}
\label{fw8i}
\lhead{:mathe:ga:}

Sind die Vektoren $e_1, e_2$ und $e_3$ die Standardbasisvektoren des
dreidimensionalen Raumens $\R^3$. So gibt es acht Basisvektoren des
$\mathscr{G}^3$.

\[
\begin{matrix}
    1 \\ 
    e_1\hspace{3ex} e_2\hspace{3ex} e_3 \\ 
    e_1e_2\hspace{3ex} e_3e_1\hspace{3ex} e_2e_3 \\
    e_1e_2e_3
\end{matrix}
\hspace{3ex}
\begin{matrix*}[l]
    \textrm{basis für }0\textrm{-Vektoren (Skalare)} \\ 
    \textrm{basis für }1\textrm{-Vektoren (Vektoren)} \\ 
    \textrm{basis für }2\textrm{-Vektoren (Bivektoren)} \\ 
    \textrm{basis für }3\textrm{-Vektoren (Trivektoren)} \\ 
\end{matrix*}
\]

Zum Beispiel mit $e_1e_2e_3$ wird das geometische Produkt der Vektoren $e_1, e_2$
und $e_3$ gemeint.

\subsection{Referenzen}
\begin{itemize}
    \item \href{f35d.pdf}{Geometrische Algebra} f35d.tex
    \item \href{81js.pdf}{Geometrisches Produkt} 81js.tex
\end{itemize}

- [$k$-Blades](kikd.md)
- [$k$-Vektoren](93t3.md)
- [Ein Multivektor](d1fv.md)

:mathe:ga:

\subsection{Literatur}
\begin{itemize}
    \item Alan Macdonald - Linear and Geometric Algebra (2021) Seite 82
\end{itemize}
\end{document}
